\section*{Capítulo \ref{ch:b3:genomics} --- Genômica e sequenciamento}

\begin{solution}{exo:b3:genomics:1}
Um \omterm{def:b3:genomics:sanger}{didesoxinucleotídeo} não tem a hidroxila 3$'$, de modo que nenhuma
ligação fosfodiéster pode ser feita com o nucleotídeo seguinte e a
cadeia para. Só com formas didesoxi, toda cadeia pararia na primeira
posição; só com formas normais, nenhuma pararia. A mistura torna a
terminação um evento aleatório em cada posição, de modo que os
produtos formam uma escada completa, um comprimento para cada base do
molde.
\end{solution}

\begin{solution}{exo:b3:genomics:2}
$4\times 10^{8}\times \qty{300}{bp} = 1.2\times 10^{11}$ bases. Humano:
$1.2\times 10^{11}/3.2\times 10^{9} \approx 38\times$. Bactéria:
$1.2\times 10^{11}/5\times 10^{6} = \num{24000}\times$.
\end{solution}

\begin{solution}{exo:b3:genomics:3}
Um contig é uma sequência contígua montada a partir de \omterm{def:b3:genomics:ngs}{leituras}
sobrepostas; um \omterm{def:b3:genomics:assembly}{scaffold} é um conjunto ordenado e orientado de contigs
com lacunas de tamanho estimado entre eles; o \omterm{def:b3:genomics:assembly}{N50} é o comprimento de
contig em que os contigs ordenados alcançam metade das bases montadas.
Aqui só os dez contigs de \qty{8}{Mb} somam \qty{80}{Mb}, além do
ponto médio de \qty{50}{Mb} (o sétimo chega a \qty{56}{Mb}):
\omterm{def:b3:genomics:assembly}{N50} $= \qty{8}{Mb}$.
\end{solution}

\begin{solution}{exo:b3:genomics:4}
O \omterm{prop:b3:genomics:sizes}{tamanho do genoma} não acompanha a complexidade do organismo: uma
cebola tem cinco vezes o DNA de um humano, um peixe pulmonado quarenta
vezes; a levedura e o verme diferem dez vezes em DNA com números de
genes semelhantes. O excesso é não codificante: elementos transponíveis
e seus restos, íntrons, repetições satélites.
\end{solution}

\begin{solution}{exo:b3:genomics:5}
$e^{-c} = 10^{-6}$ dá $c = 6\ln 10 = 13.8$. Para $c = 8$: $N =
cG/L = 8\times 5\times 10^{6}/150 \approx \num{267000}$ \omterm{def:b3:genomics:ngs}{leituras},
$\theta
= 30/150 = 0.2$, contigs $= N e^{-6.4} = \num{267000}\times 0.00166
\approx 440$.
\end{solution}

\begin{solution}{exo:b3:genomics:6}
$P(X < 8) = P(X \le 7) \approx P\bigl(Z < (7.5 - 15)/2.74\bigr) =
P(Z < -2.74) \approx 0.003$. A $30\times$ os dois alelos de um
heterozigoto são quase sempre vistos muitas vezes, a cobertura é
desigual (regiões ricas em GC recebem menos \omterm{def:b3:genomics:ngs}{leituras}, de modo que
alguns sítios ficam com metade da média) e restam \omterm{def:b3:genomics:ngs}{leituras} bastantes
para separar um alelo verdadeiro dos erros de $10^{-3}$.
\end{solution}

\begin{solution}{exo:b3:genomics:7}
Toda \omterm{def:b3:genomics:ngs}{leitura} de \qty{150}{bp} vinda de dentro da repetição é idêntica
venha da cópia que vier, de modo que o montador constrói um único nó
de \qty{6}{kb} com $500$ caminhos entrando e $500$ saindo; ele não tem
como dizer qual entrada casa com qual saída, e a \omterm{def:b3:genomics:assembly}{montagem} se parte em
$500$ lacunas, com a repetição presente uma única vez e com cobertura
$500$ vezes a média. \omterm{def:b3:genomics:ngs}{Leituras} mais longas que a repetição mais flancos
únicos dos dois lados --- \qty{8}{kb} ou mais --- atravessam cada cópia
e a resolvem.
\end{solution}

\begin{solution}{exo:b3:genomics:8}
$4.8\times 10^{7}/1000 = \num{48000}$ diferenças codificadoras; um
terço, cerca de \num{16000}, altera a proteína.
\end{solution}

\begin{solution}{exo:b3:genomics:9}
Falsos positivos esperados $10^{6}\times 5\times 10^{-8} = 0.05$. Um
risco relativo de $1.15$ sobre uma doença de \qty{2}{\%} muda alguns
casos por mil; mil pessoas contêm cerca de vinte casos, longe demais
do bastante para vê-lo (o poder cresce com o número de casos e com o
quadrado do efeito). Para um indivíduo, $0.3$ ponto percentual não
significa nada. Mas o alelo marca um gene cuja mudança modesta de
atividade altera o risco da doença --- o gene, e sua via, podem ser um
alvo farmacológico cuja inibição completa tenha grande efeito.
\end{solution}

\begin{solution}{exo:b3:genomics:10}
De genética de populações: nas bactérias, com populações enormes, a
seleção remove até inserções levemente custosas; nos mamíferos, com
populações efetivas pequenas, a deriva deixa acumular DNA não
codificante levemente deletério --- apoiada pela correlação inversa
entre \omterm{prop:b3:genomics:sizes}{tamanho do genoma} e tamanho populacional ao longo das linhagens,
enfraquecida pelas exceções. Estrutural: os genomas eucariontes são
invadidos por transpósons que as bactérias, com um viés a favor da
deleção e sem refúgio meiótico para elementos egoístas, expurgam ---
apoiada pela correlação do \omterm{prop:b3:genomics:sizes}{tamanho do genoma} com o conteúdo de
transpósons. Regulatória: um desenvolvimento complexo precisa de mais
DNA regulatório --- apoiada pelos \omterm{def:b3:genomics:comparative}{elementos não codificantes
conservados} em torno dos genes do desenvolvimento, enfraquecida pelo
fato de que só \qty{8}{\%} do genoma humano mostra alguma seleção, de
modo que a maior parte do DNA não codificante não é regulatória.
\end{solution}

\begin{solution}{exo:b3:genomics:11}
Os inícios dos dois tipos caem com densidade total $\rho = (N_{1} +
N_{2})/G$. Uma \omterm{def:b3:genomics:ngs}{leitura} de comprimento $L_{i}$ é a mais à direita de seu
contig se nenhuma \omterm{def:b3:genomics:ngs}{leitura} de qualquer dos tipos começa nas $L_{i} - T$
bases após ela: probabilidade $e^{-\rho(L_{i} - T)}$. Logo, contigs
$= N_{1}
e^{-\rho(L_{1} - T)} + N_{2} e^{-\rho(L_{2} - T)}$. Uma
\omterm{def:b3:genomics:ngs}{leitura} longa é a mais à direita com probabilidade exponencialmente
menor que uma curta, e cada \omterm{def:b3:genomics:ngs}{leitura} longa também elimina a chance de
uma lacuna ao longo de todo o seu comprimento; o mesmo número de bases
em \omterm{def:b3:genomics:ngs}{leituras} curtas acrescenta muitos pontos de início, mas cada janela
protegida é curta, e por isso as \omterm{def:b3:genomics:ngs}{leituras} longas vencem.
\end{solution}

\begin{solution}{exo:b3:genomics:12}
Duas \omterm{def:b3:genomics:comparative}{duplicações do genoma inteiro} preveem que o padrão de quatro
cópias seja genômico: quartetos de segmentos cromossômicos \omterm{def:b3:genomics:comparative}{parálogos}
(paralogons) carregando as mesmas famílias gênicas na mesma ordem, com
as duplicações datadas do mesmo momento para todas as famílias; um
único agrupamento no anfioxo e em \emph{Ciona}, que divergiram antes
dos eventos; e nas lampreias, que divergiram por volta deles, um
estado intermediário ou derivado de forma independente. Duplicações
independentes só do agrupamento Hox preveem paralogons apenas para o
Hox, vizinhos com histórias diferentes e datas de duplicação
diferentes entre as famílias. Os paralogons genômicos e a datação
compartilhada, como se observa, favorecem a \omterm{def:b3:genomics:comparative}{duplicação do genoma
inteiro}.
\end{solution}

\begin{solution}{pb:b3:genomics:1}
\textbf{1.} $0.002\times 6\times 10^{8} = 1.2\times 10^{6}$ \omterm{def:b3:genomics:ngs}{leituras};
$c = 1.2\times 10^{6}\times 150/4.6\times 10^{6} \approx 39$.
\textbf{2.} $e^{-39} \approx 10^{-17}$: na prática, nenhuma base fica
descoberta.
\textbf{3.} $\theta = 0.2$; contigs $= 1.2\times 10^{6} e^{-31} \approx
0$: um contig em teoria, com lacunas apenas nas repetições.
\textbf{4.} $c = \num{30000}\times 150/4.6\times 10^{6} = 0.98$; fração
não coberta $e^{-0.98} = 0.38$, cerca de \qty{1.7}{Mb}; contigs $= \num{30000}
\times e^{-0.78} \approx \num{13700}$.
\textbf{5.} Os sete óperons dão \omterm{def:b3:genomics:ngs}{leituras} idênticas e colapsam num único
contig de \qty{5}{kb}, ao qual chegam sete flancos esquerdos únicos e do
qual saem sete flancos direitos cujo pareamento é desconhecido: $14$
pontas de contig, sete lacunas.
\textbf{6.} Cerca de \num{4600} genes; sequência codificadora
$0.88\times 4.6\times 10^{6}
= 4.0\times 10^{6}$ bp, gene médio de
\qty{880}{bp}.
\textbf{7.} $0.998\times 6\times 10^{8}\times 150/3.2\times 10^{9}
\approx 28\times$.
\textbf{8.} Cerca de $14$ \omterm{def:b3:genomics:ngs}{leituras} por alelo.
\textbf{9.} $28\times 10^{-3}/3 \approx 0.009$ \omterm{def:b3:genomics:ngs}{leitura} mostrando uma
dada base errada --- a chance de três ou mais é de cerca de $10^{-7}$
--- e \qty{20}{\%} de $28$ são $5.6$ \omterm{def:b3:genomics:ngs}{leituras}; um erro não passa em
nenhum dos dois testes, ao passo que um alelo heterozigoto verdadeiro é
esperado em $14$.
\textbf{10.} $3.2\times 10^{9}/1500 \approx 2.1\times 10^{6}$ sítios
heterozigotos; em sequência codificadora, $4.8\times 10^{7}/1500 \approx
\num{32000}$.
\textbf{11.} Uma \omterm{def:b3:genomics:ngs}{leitura} de \qty{150}{bp} vinda de um Alu casa quase
igualmente bem com muitas das $1.1$ milhão de cópias, e por isso é
colocada na cópia errada ou recebe baixa confiança de mapeamento. As
diferenças entre as cópias passam então a parecer variantes
heterozigotas, e as variantes verdadeiras ficam escondidas entre elas:
as \omterm{met:b3:genomics:pipeline}{chamadas de variantes} dentro de Alu costumam ser filtradas, e os
elementos são pontos cegos do sequenciamento de \omterm{def:b3:genomics:ngs}{leituras} curtas.
\textbf{12.} $1.2\times 10^{-8}\times 6.4\times 10^{9} \approx 77$
mutações novas. Cerca de $14$ \omterm{def:b3:genomics:ngs}{leituras} deveriam mostrar a variante na
criança; mas, se o sítio de um dos pais é coberto por apenas cinco
\omterm{def:b3:genomics:ngs}{leituras}, um alelo heterozigoto passa despercebido com probabilidade
$2^{-5} = 3\,\%$ e uma variante herdada é chamada de nova --- por isso
os pais precisam ser sequenciados tão profundamente quanto a criança.
\textbf{13.} $4.8\times 10^{7}\times 100/150 = 3.2\times 10^{7}$
\omterm{def:b3:genomics:ngs}{leituras}, vinte vezes menos que o genoma: o custo de sequenciamento cai
vinte vezes, mais do que o custo da etapa de captura.
\textbf{14.} $1.1\times 10^{6}\times 300 = 3.3\times 10^{8}$ bp, cerca
de \qty{10}{\%} do genoma; \qty{10}{\%} das \omterm{def:b3:genomics:ngs}{leituras}, uns $6\times
10^{7}$.
\textbf{15.} $\num{60000}\times 1500 = 9\times 10^{7}$ bp; $9\times
10^{7}/1.6\times 10^{10} = 0.56\,\%$.
\textbf{16.} $0.012\times 2.9\times 10^{9} = 3.5\times 10^{7}$
substituições, $1.7\times 10^{7}$ por linhagem; taxa $1.7\times
10^{7}/(2.9\times 10^{9}\times 6.5\times 10^{6}) = 9\times 10^{-10}$
por base por ano, $2.3\times 10^{-8}$ por geração de 25 anos --- cerca
do dobro da taxa medida em genealogias, $1.2\times 10^{-8}$, uma
discrepância conhecida que sugere gerações mais longas no passado ou
uma separação mais antiga.
\textbf{17.} Quatro cópias de \num{16000} dariam \num{64000};
\num{20000} permanecem, de modo que, das \num{48000} cópias extras, só
\num{4000} sobrevivem: \qty{92}{\%} das duplicatas se perderam.
\textbf{18.} O verme e o humano têm o mesmo número de genes; a
complexidade está em como eles são usados --- o splicing alternativo
multiplica um gene em milhares de proteínas (\num{38016} no
\emph{Dscam}), a regulação combina fatores de transcrição e
intensificadores de maneiras específicas de cada tipo celular, os \omterm{def:b3:rna-regulation:ncrna}{RNAs
não codificantes} acrescentam camadas, e as proteínas interagem em
redes.
\textbf{19.} Elementos regulatórios (intensificadores, promotores,
isolantes), genes de \omterm{def:b3:rna-regulation:ncrna}{RNA não codificante}, sinais de splicing e de
localização, origens e sequências estruturais. Teste um candidato a
intensificador colocando-o antes de um gene repórter num embrião
transgênico e vendo onde o repórter é expresso, depois deletando o
elemento no genoma com CRISPR e medindo os genes vizinhos.
\textbf{20.} $0.05/10^{6} = 5\times 10^{-8}$; a $p < 10^{-5}$,
$10^{6}
\times 10^{-5} = 10$ falsos positivos esperados.
\textbf{21.} Chances do não portador $0.009/0.991 = 0.00908$; com uma
cópia, chances $0.00954$, risco $0.945\,\%$; com duas cópias, chances
$0.01001$, risco $0.99\,\%$.
\textbf{22.} O escore soma, sobre os $200$ alelos, o número de cópias de
risco que a pessoa carrega, ponderado pelo logaritmo da razão de
chances de cada alelo. A contagem tem média $200\times 2\times 0.3 = 120$ e desvio padrão $\sqrt{200\times 2\times 0.3\times 0.7} \approx 9$;
o \qty{1}{\%} mais alto carrega cerca de $21$ alelos acima da média, e
$1.05^{21} \approx
2.8$: quase o triplo das chances de uma pessoa
média, a partir de alelos que individualmente mudam o risco em
\qty{5}{\%}.
\textbf{23.} Os alelos dentro de um bloco de dezenas de quilobases são
herdados juntos (desequilíbrio de ligação), de modo que qualquer um
deles mostra a mesma associação; o SNP genotipado é apenas o que estava
no arranjo. Para achar a variante causal: sequencie a região em casos e
controles, estreite o conjunto estatisticamente e depois teste cada
candidata --- ensaios com repórter para a atividade intensificadora de
cada alelo, edição de cada alelo em células com medida da expressão dos
genes vizinhos, colocalização com sinais de loco de característica
quantitativa de expressão.
\textbf{24.} $1 - e^{-0.5} = 39\,\%$ do genoma é lido. Ele dá uma
amostra direta de uma população num momento conhecido do passado ---
frequências alélicas antes de migrações e mesclas posteriores, o
comprimento dos segmentos neandertais (longos, porque poucas gerações
de recombinação os haviam quebrado) e daí uma data para a mescla ---
coisas que os genomas modernos só conseguem inferir por modelos.
\textbf{25.} Cerca de \num{13700} contigs no piloto; $28\times$ por
genoma haploide, cerca de $14$ \omterm{def:b3:genomics:ngs}{leituras} por alelo; umas $77$ mutações
novas numa criança.
\end{solution}
